\section{启发式惯性评分，而不是已证定律}

\begin{definition}[启发式惯性评分]
如果分析者为 $\Pi(s)$ 的各个分量选择了共同的归一化方式，就可以引入一个\textbf{启发式惯性评分}：
\[
  \widehat{\FI}(s) = \alpha \Reward(s) + \beta \Switch(s) + \gamma \Env(s),
\]
其中 $\alpha,\beta,\gamma > 0$ 是建模权重。
\end{definition}

\begin{remark}
$\widehat{\FI}$ 上方的帽号很重要。这个表达式是一个记账装置，而不是某条已经成立的行为动力学定律。当需要一个单一比较分数时，它很有用；但手稿不应假装这些权重已经被实证识别，也不应假装线性形式是唯一正确的。
\end{remark}

\begin{discussion}
这里真正站得住脚的，是\textbf{支架的结构}，而不是“这个线性组合已经被验证为普适公式”的说法。教材与综述文献为使用低维状态描述和稳定性语言提供了正当性。它们并不支持我们宣称，人类行动成本服从某条已知的三项线性定律。这个评分之所以被保留，是因为它在操作层面上有帮助；但它被明确从“形式定律”降级成了“启发式模型”。
\end{discussion}

\begin{remark}
本章原本也可以完全避免标量评分，而仅保持序关系。之所以保留 $\widehat{\FI}$，是出于教学上的经济性：后续章节需要一种紧凑方式来表达一个状态比另一个状态更难撼动。因此，本章批准它作为一种\emph{比较辅助}，而不是某个深层实证结果。
\end{remark}